



\hyphenpenalty=1000 %permet d'empêcher la coupure des mots avec des tirets.



%\usepackage{showframe}



%-------------------------------------------------------------------------------------------------------------------------------------------------------------------------------------------------------------
%---- La commande ci-dessous permet de modifier la largeur de l'en-tête et du pied de page afin de couvrir le titre de sections. Il est tiré d'un exemple de la documentation du package "fancyhdr" (page 13).
%-------------------------------------------------------------------------------------------------------------------------------------------------------------------------------------------------------------
\fancyhfoffset[l]{0.2cm} % le "l" en paramètre permet d'indiquer qu'on ne veut modifier que la marge à gauche.
\renewcommand{\headrule}{{%
\hrule \headwidth \headrulewidth \vskip-\headrulewidth}}
\renewcommand\footrulewidth{\headrulewidth}
\renewcommand{\footrule}{{%
\vskip-\footruleskip\vskip-\footrulewidth
\hrule \headwidth \footrulewidth\vskip\footruleskip}}

%-----------------------------------------------------------------
%---- Modification présentation de la page: marges de la page ----
%-----------------------------------------------------------------
%\addtolength{\hoffset}{-1in}              % 1
%\addtolength{\voffset}{-1in}              % 2
\addtolength{\oddsidemargin}{-0.4 in} % 3
\addtolength{\evensidemargin}{-1.4in} % 3
\addtolength{\topmargin}{-1in}       % 4
\addtolength{\headheight}{6.1pt}       % 5
%\addtolength{\headsep}{-0.2cm}           % 6
\setlength{\textheight}{26cm}    % 7
\setlength{\textwidth}{17.5cm}      % 8
\addtolength{\marginparsep}{0pt}      % 9
\setlength{\marginparwidth}{0pt}   % 10
\addtolength{\footskip}{-1mm}           %11












%%%%%%%%%%%%%%%%%%%%%%%%%%%%%%%%%%%%%%%%%%%%%%%%%%%%%%%%%%%%%%%%%%%
%%%%%%%%%%%%%%%%%%%%%%%%%%%   POLICES   %%%%%%%%%%%%%%%%%%%%%%%%%%%
%%%%%%%%%%%%%%%%%%%%%%%%%%%%%%%%%%%%%%%%%%%%%%%%%%%%%%%%%%%%%%%%%%%
%-----------------------------------------------------------------------------------------------------------------------------------------
%---- ETABLISSEMENT DE LA POLICE "Computer Modern" PAR DEFAUT. En effet, en utilisant XeLaTeX la police par défaut est "Latin Modern" ----
%-----------------------------------------------------------------------------------------------------------------------------------------
%\renewcommand\rmdefault{cmr}
%\renewcommand\sfdefault{cmss}
%\renewcommand\ttdefault{cmtt}
%
%%-------------------------------------------------------------------------
%%---- Polices chapitre, sections, ... ----
%%-------------------------------------------------------------------------
%\setcounter{chapter}{0}
%\shadowoffset{1pt}
\chapterfont{\raggedleft\vspace{-0.5cm}\fontsize{36}{28}\fontshape{sl}\selectfont\shadowtext}
\chaptertitlefont{\raggedleft\fontsize{25}{20}\fontshape{sl}\selectfont\shadowtext}
%\shadowoffset{0.7pt}
\sectionfont{\vspace{0.6cm}\hspace{-15mm}\fontsize{17}{10}\fontshape{sl}\selectfont }
%%\sectionfont{\shadowtext\ulemheading{\uuline}\selectfont}
%\subsectionfont{\vspace{0.3cm}\hspace{-10mm}\fontsize{14}{10}\fontshape{sl}\selectfont\vspace{0mm} }


%\titleformat{\section}
%  {\normalfont\LARGE\bfseries}
%  {}
%  {0em}
%  {\thesection\hspace{3mm}\underline}


%\titleformat
%{\section} % command
%[hang] % shape
%{\bfseries\Large\itshape} % format
%{\thesection} % label
%{0mm} % sep
%{
%} % before-code
%[
%\vspace{-0.5ex}%
%] % after-code


%------------------------------------------
%---- Déclaration de nouvelles polices ----
%------------------------------------------
%\DeclareTextFontCommand{\poleclairage}{\fontfamily{comicneue}\selectfont}
\DeclareFixedFont{\poleclairage}{OT1}{pcr}{m}{n}{12}
% \DeclareFixedFont {?cmd?} {?encoding?} {?family?} {?series?} {?shape?} {?size?}
%\fontencoding {encoding}
%\fontfamily {family}
%\fontseries {series}
%\fontshape {shape}
%\fontsize {size} {baselineskip}
%\linespread {factor}














\frenchbsetup{StandardLists=true} % Chez les anglo-saxons (créateur de LATEX) les listes ont des puces ronde () et en français des tirets (–). L’option frenchb du package babel se charge habituellement de la transformation. Le rôle du package enumitem est également de modifier les puces des listes. Deux packages qui font la même chose ne cohabitent par forcement bien. La solution (si on charge enumitem) est de demander à frenchb de ne pas s’occuper des listes, par la commande \frenchbsetup{StandardLists=true}



%\raggedbottom  % Empeche le texte de s'étaler verticalement sur une page (-> par défaut avec l'option "twoside")




\newcommand{\cen}[1]{\begin{center}
#1\vspace*{-4mm}\end{center}}


\newcommand{\Separationitems}{3mm} % Espace vertical entre les items
\newcommand{\margegauche}{7mm} % Espace horizontal après le numéro de l'exercice
\newcommand{\EspaceApresItems}{3mm}
\SetEnumitemKey{deuxcol}{
	itemsep = \Separationitems,		% Espace vertical entre les items
	leftmargin = \margegauche,		% Espace horizontal après le numéro de l'exercice
	labelsep*= \EspaceApresItems,
	label  = \alph*),
	before  = \raggedcolumns \begin{multicols}{2},
	after   = \end{multicols}}

\SetEnumitemKey{troiscol}{
	itemsep = \Separationitems,
	leftmargin = \margegauche,
	labelsep*= \EspaceApresItems,
	label  = \alph*),
	before  = \raggedcolumns \begin{multicols}{3},
	after   = \end{multicols}}
	
\SetEnumitemKey{quatrecol}{
	itemsep = \Separationitems,
	leftmargin = \margegauche,
	labelsep*= \EspaceApresItems,
	leftmargin = *,
	label  = \alph*),
	before  = \raggedcolumns \begin{multicols}{4},
	after   = \end{multicols}}

\SetEnumitemKey{cinqcol}{
	itemsep = \Separationitems,
%	leftmargin = \margegauche,
	labelsep*= \EspaceApresItems,
	label  = \alph*),
	before  = \raggedcolumns \begin{multicols}{5},
	after   = \end{multicols}}

\SetEnumitemKey{sixcol}{
	itemsep = \Separationitems,
	leftmargin = \margegauche,
	labelsep*= \EspaceApresItems,
	label  = \alph*),
	before  = \raggedcolumns \begin{multicols}{6},
	after   = \end{multicols}}

\SetEnumitemKey{septcol}{
	itemsep = \Separationitems,
	leftmargin = \margegauche,
	labelsep*= \EspaceApresItems,
	label  = \alph*),
	before  = \raggedcolumns \begin{multicols}{6},
	after   = \end{multicols}}
	
\SetEnumitemKey{huitcol}{
	itemsep = \Separationitems,
	leftmargin = \margegauche,
	labelsep*= \EspaceApresItems,
	label  = \alph*),
	before  = \raggedcolumns \begin{multicols}{6},
	after   = \end{multicols}}


% Ici on défini l'espace après les items de niveau 2 (sous-items) POUR TOUTES LES LISTES : (->définit localement dans le fichier maître)
%\setlist[2]{labelsep*=3mm}



